\documentclass[11pt]{article}

\usepackage[utf8]{inputenc}
\usepackage[T1]{fontenc}
\usepackage{lmodern}
\usepackage[margin=1in]{geometry}
\usepackage{amsmath, amssymb, amsfonts, mathtools}
\usepackage{booktabs}
\usepackage{array}
\usepackage{enumitem}
\usepackage{hyperref}

\hypersetup{
  colorlinks=true,
  linkcolor=blue,
  urlcolor=blue,
  pdftitle={Kamada-Kawai and Fruchterman-Reingold Graph Layout Methods},
  pdfauthor={Codex}
}

\setlength{\parskip}{0.6em}
\setlength{\parindent}{0pt}

\newcommand{\R}{\mathbb{R}}
\newcommand{\vect}[1]{\boldsymbol{#1}}
\newcommand{\norm}[1]{\left\lVert #1 \right\rVert}

\title{Kamada--Kawai and Fruchterman--Reingold Graph Layout Methods\\
\large A Mathematical Comparison of Two Classical Force-Directed Layout Algorithms}
\author{Generated for the \texttt{grip} project}
\date{\today}

\begin{document}

\maketitle

\begin{abstract}
This document explains in detail two classical force-directed graph layout
methods: the Kamada--Kawai (KK) algorithm and the Fruchterman--Reingold (FR)
algorithm. For each method, we describe the mathematical model, define the
meaning of the main terms and parameters, write the update equations, and
explain the optimization logic. We then compare the two methods in terms of
their objectives, force laws, computational behavior, and practical strengths
and weaknesses.
\end{abstract}

\section{Preliminaries and notation}

Let
\[
G = (V, E)
\]
be an undirected graph with vertex set $V = \{1, 2, \dots, n\}$ and edge set
$E$. A layout in two dimensions assigns each vertex $i \in V$ a position
\[
\vect{x}_i = (x_i, y_i)^\top \in \R^2.
\]
Collectively, the layout is
\[
X = (\vect{x}_1, \dots, \vect{x}_n).
\]

For any pair of vertices $i$ and $j$, define:
\begin{itemize}[leftmargin=1.5em]
\item \textbf{Euclidean distance in the drawing}
\[
r_{ij}(X) = \norm{\vect{x}_i - \vect{x}_j}.
\]
\item \textbf{Graph-theoretic distance}
\[
g_{ij},
\]
the length of a shortest path between $i$ and $j$ in the graph. In an
unweighted graph this is the number of edges on a shortest path; in a weighted
graph it can be the shortest weighted path length.
\item \textbf{Graph diameter}
\[
D = \max_{i,j \in V} g_{ij},
\]
assuming the graph is connected.
\end{itemize}

Both KK and FR are \emph{force-directed} methods: they attempt to place
vertices by balancing attractive and repulsive tendencies so that the final
drawing reflects graph structure. The important difference is that KK is built
around matching graph distances, whereas FR is built around balancing local
forces to create a visually even drawing.

\subsection{Meaning of the most common terms}

\begin{center}
\begin{tabular}{>{\raggedright\arraybackslash}p{0.22\textwidth} >{\raggedright\arraybackslash}p{0.68\textwidth}}
\toprule
\textbf{Term} & \textbf{Meaning} \\
\midrule
Layout & The map from each vertex to a point in the plane. \\
Euclidean distance $r_{ij}$ & The actual geometric distance between two vertices in the current drawing. \\
Graph distance $g_{ij}$ & The shortest-path distance between two vertices in the graph itself. \\
Ideal distance & A desired geometric separation between two vertices. In KK this depends on $g_{ij}$. In FR the global spacing parameter $k$ plays an analogous role. \\
Attractive force & A force pulling adjacent or structurally related vertices closer. \\
Repulsive force & A force pushing vertices apart to avoid crowding and overlap. \\
Energy / objective & A scalar function the algorithm tries to reduce. KK has an explicit energy; FR is usually presented as an iterative force simulation rather than a single explicit objective. \\
Temperature & A cap on how far a vertex may move in one FR iteration; it stabilizes the algorithm and gradually decreases. \\
Convergence & The state in which updates become very small or a stopping rule is reached. \\
\bottomrule
\end{tabular}
\end{center}

\section{The Kamada--Kawai method}

\subsection{Core idea}

The Kamada--Kawai method treats the graph as a complete system of springs.
Every pair of vertices $(i,j)$ is connected by a virtual spring, not only the
edges in $E$. The desired spring length is proportional to the graph distance
$g_{ij}$. Therefore:
\begin{itemize}[leftmargin=1.5em]
\item vertices that are close in the graph should be close in the picture,
\item vertices that are far apart in the graph should be far apart in the picture.
\end{itemize}

This makes KK a graph-distance-preserving method. It is closely related to
weighted stress minimization.

\subsection{Ideal lengths and spring strengths}

KK defines an ideal geometric length $\ell_{ij}$ and a spring constant
$k_{ij}$ for every unordered pair $i \neq j$:
\[
\ell_{ij} = L_0 \, g_{ij}, \qquad L_0 = \frac{L}{D},
\]
and
\[
k_{ij} = \frac{K}{g_{ij}^2}.
\]

Here:
\begin{itemize}[leftmargin=1.5em]
\item $L$ is a user-chosen global scale for the drawing,
\item $D$ is the graph diameter,
\item $L_0$ converts graph distances into drawing-space distances,
\item $K > 0$ is a global stiffness scale,
\item $\ell_{ij}$ is the desired geometric distance between $i$ and $j$,
\item $k_{ij}$ is the spring stiffness between $i$ and $j$.
\end{itemize}

The formula $k_{ij} = K / g_{ij}^2$ means short graph distances get stronger
springs and large graph distances get weaker springs. Intuitively, nearby graph
relationships are enforced more strongly than distant ones.

\subsection{Energy function}

KK defines the total energy
\[
E(X) = \frac{1}{2} \sum_{1 \le i < j \le n}
k_{ij} \left(r_{ij}(X) - \ell_{ij}\right)^2.
\]

This formula is central. Each term has a precise meaning:
\begin{itemize}[leftmargin=1.5em]
\item $r_{ij}(X)$ is the current geometric distance in the layout,
\item $\ell_{ij}$ is the desired distance implied by the graph,
\item $r_{ij}(X) - \ell_{ij}$ is the distance error for that pair,
\item $\left(r_{ij}(X) - \ell_{ij}\right)^2$ penalizes large errors more heavily,
\item $k_{ij}$ weights that penalty by the importance of the pair,
\item the sum adds these penalties across all pairs of vertices.
\end{itemize}

Therefore KK tries to make Euclidean distances match scaled graph distances.
If $r_{ij}(X) = \ell_{ij}$ for every pair, then the energy is zero. In practice
that is usually impossible, so the algorithm seeks a low-energy compromise.

\subsection{Gradient with respect to one vertex}

KK optimizes one vertex at a time. Fix a vertex $m$. Its contribution to the
first derivatives is
\[
\frac{\partial E}{\partial x_m}
= \sum_{\substack{i=1 \\ i \ne m}}^n
k_{mi}
\left[
(x_m - x_i)
- \ell_{mi}\frac{x_m - x_i}{r_{mi}}
\right],
\]
\[
\frac{\partial E}{\partial y_m}
= \sum_{\substack{i=1 \\ i \ne m}}^n
k_{mi}
\left[
(y_m - y_i)
- \ell_{mi}\frac{y_m - y_i}{r_{mi}}
\right],
\]
where
\[
r_{mi} = \sqrt{(x_m - x_i)^2 + (y_m - y_i)^2}.
\]

It is useful to define the local residual magnitude
\[
\Delta_m
=
\sqrt{
\left(\frac{\partial E}{\partial x_m}\right)^2
+
\left(\frac{\partial E}{\partial y_m}\right)^2
}.
\]

Interpretation:
\begin{itemize}[leftmargin=1.5em]
\item $\Delta_m$ measures how far vertex $m$ is from local equilibrium,
\item a large $\Delta_m$ means the forces around $m$ are still badly unbalanced,
\item KK usually selects the vertex with the largest $\Delta_m$ and optimizes it next.
\end{itemize}

\subsection{Second derivatives and Newton step}

KK uses a Newton--Raphson update for the selected vertex $m$. The relevant
second derivatives are
\[
\frac{\partial^2 E}{\partial x_m^2}
= \sum_{\substack{i=1 \\ i \ne m}}^n
k_{mi}
\left[
1 - \ell_{mi}\frac{(y_m - y_i)^2}{r_{mi}^3}
\right],
\]
\[
\frac{\partial^2 E}{\partial y_m^2}
= \sum_{\substack{i=1 \\ i \ne m}}^n
k_{mi}
\left[
1 - \ell_{mi}\frac{(x_m - x_i)^2}{r_{mi}^3}
\right],
\]
\[
\frac{\partial^2 E}{\partial x_m \partial y_m}
= \sum_{\substack{i=1 \\ i \ne m}}^n
k_{mi}
\ell_{mi}
\frac{(x_m - x_i)(y_m - y_i)}{r_{mi}^3}.
\]

The Newton step $(\delta x_m, \delta y_m)$ solves
\[
\begin{bmatrix}
\dfrac{\partial^2 E}{\partial x_m^2} &
\dfrac{\partial^2 E}{\partial x_m \partial y_m} \\[1.2ex]
\dfrac{\partial^2 E}{\partial x_m \partial y_m} &
\dfrac{\partial^2 E}{\partial y_m^2}
\end{bmatrix}
\begin{bmatrix}
\delta x_m \\ \delta y_m
\end{bmatrix}
=
-
\begin{bmatrix}
\dfrac{\partial E}{\partial x_m} \\[1.2ex]
\dfrac{\partial E}{\partial y_m}
\end{bmatrix}.
\]

Then the vertex is updated by
\[
x_m \leftarrow x_m + \delta x_m,
\qquad
y_m \leftarrow y_m + \delta y_m.
\]

This is repeated until $\Delta_m$ becomes sufficiently small; then the
algorithm chooses the next vertex with largest residual and continues.

\subsection{Algorithmic outline}

A standard KK workflow is:
\begin{enumerate}[leftmargin=1.5em]
\item Compute all-pairs shortest-path distances $g_{ij}$.
\item Define $\ell_{ij}$ and $k_{ij}$ from $g_{ij}$.
\item Start from an initial placement.
\item Compute $\Delta_m$ for every vertex.
\item Choose the vertex with largest $\Delta_m$.
\item Apply Newton updates to that vertex until its local residual is small.
\item Repeat until the largest residual in the graph is below a threshold.
\end{enumerate}

\subsection{Why the method works}

KK works because it turns graph drawing into a continuous optimization problem.
The graph tells us what pairwise distances \emph{should} be; the optimization
tries to realize those distances geometrically. The method is especially good
when preserving global graph distances matters.

\subsection{Strengths and limitations}

\textbf{Strengths}
\begin{itemize}[leftmargin=1.5em]
\item It has a clear mathematical objective.
\item It often produces very good layouts for small and medium-sized graphs.
\item It is especially strong on graphs where preserving shortest-path
structure is visually meaningful.
\end{itemize}

\textbf{Limitations}
\begin{itemize}[leftmargin=1.5em]
\item It requires all-pairs shortest-path distances.
\item The objective involves all pairs of vertices, so time and memory can grow quickly.
\item It can become slow on large graphs.
\item Like many nonconvex optimization methods, it can get trapped in local minima.
\end{itemize}

\subsection{Computational cost}

If implemented naively:
\begin{itemize}[leftmargin=1.5em]
\item all-pairs shortest paths cost at least quadratic and often cubic time,
\item the energy and derivatives involve all pairs, leading to roughly
$O(n^2)$ work per global sweep,
\item memory can also become $O(n^2)$ if all pairwise quantities are stored.
\end{itemize}

This is the main reason KK is often reserved for smaller graphs or used inside
multilevel or approximate schemes.

\section{The Fruchterman--Reingold method}

\subsection{Core idea}

The Fruchterman--Reingold method models a graph drawing using two competing
forces:
\begin{itemize}[leftmargin=1.5em]
\item \textbf{Repulsion} between every pair of vertices, which spreads the graph out,
\item \textbf{Attraction} along graph edges, which pulls adjacent vertices together.
\end{itemize}

Instead of matching graph distances directly, FR seeks a visually balanced state
in which vertices are evenly distributed while edges stay reasonably short.

\subsection{Global spacing parameter}

FR introduces a single characteristic distance
\[
k = C \sqrt{\frac{A}{n}},
\]
where:
\begin{itemize}[leftmargin=1.5em]
\item $A$ is the area available for drawing,
\item $n$ is the number of vertices,
\item $C$ is a constant, often taken as $1$ in simplified descriptions,
\item $k$ is the \emph{ideal spacing scale}.
\end{itemize}

The parameter $k$ is not a target distance for every pair in the same sense as
KK's $\ell_{ij}$. Instead, it sets the overall length scale of the drawing. If
the graph is drawn in a larger area, then $k$ grows; if more vertices are packed
into the same area, then $k$ shrinks.

\subsection{Force laws}

Let
\[
\vect{\delta}_{ij} = \vect{x}_i - \vect{x}_j,
\qquad
r_{ij} = \norm{\vect{\delta}_{ij}}.
\]

The classical FR scalar force magnitudes are
\[
f_r(r) = \frac{k^2}{r}
\qquad \text{(repulsive force)}
\]
and
\[
f_a(r) = \frac{r^2}{k}
\qquad \text{(attractive force)}.
\]

To convert these scalar magnitudes into vectors, multiply by the unit direction
\[
\frac{\vect{\delta}_{ij}}{r_{ij}}.
\]
Thus the repulsive force on vertex $i$ from vertex $j$ is
\[
\vect{F}^{\text{rep}}_{ij}
=
\frac{\vect{\delta}_{ij}}{r_{ij}} f_r(r_{ij})
=
\frac{k^2}{r_{ij}^2}\,\vect{\delta}_{ij}.
\]

For an edge $(i,j) \in E$, the attractive force on vertex $i$ due to vertex $j$
is
\[
\vect{F}^{\text{att}}_{ij}
=
-
\frac{\vect{\delta}_{ij}}{r_{ij}} f_a(r_{ij})
=
-
\frac{r_{ij}}{k}\,\vect{\delta}_{ij}.
\]

These formulas mean:
\begin{itemize}[leftmargin=1.5em]
\item very close vertices repel strongly because $f_r(r) = k^2/r$ grows as $r \to 0$,
\item long edges attract strongly because $f_a(r) = r^2/k$ grows with $r$,
\item the drawing is pushed toward a balance between spreading and coherence.
\end{itemize}

\subsection{Meaning of the terms in the FR formulas}

\begin{itemize}[leftmargin=1.5em]
\item $\vect{\delta}_{ij}$ is the displacement vector from vertex $j$ to vertex $i$.
\item $r_{ij}$ is the current Euclidean distance between the two vertices.
\item $k$ is the global ideal spacing scale.
\item $f_r(r)$ is the scalar repulsive magnitude.
\item $f_a(r)$ is the scalar attractive magnitude.
\item $\vect{F}^{\text{rep}}_{ij}$ is the vector pushing $i$ away from $j$.
\item $\vect{F}^{\text{att}}_{ij}$ is the vector pulling $i$ toward $j$ if $(i,j)$ is an edge.
\end{itemize}

An important observation is that if one informally equates the magnitudes
$f_a(r)$ and $f_r(r)$, then
\[
\frac{r^2}{k} = \frac{k^2}{r}
\quad \Longrightarrow \quad
r = k.
\]
So $k$ represents a natural spacing scale at which attraction and repulsion are
balanced in the simplest pairwise sense.

\subsection{Displacement accumulation}

At each iteration, FR accumulates a displacement vector for every vertex:
\[
\vect{d}_i =
\sum_{\substack{j=1 \\ j \ne i}}^n \vect{F}^{\text{rep}}_{ij}
+
\sum_{\substack{j:(i,j)\in E}} \vect{F}^{\text{att}}_{ij}.
\]

In practice, implementations often compute this in two stages:
\begin{enumerate}[leftmargin=1.5em]
\item add repulsive contributions for all vertex pairs,
\item add attractive contributions for all edges.
\end{enumerate}

\subsection{Temperature-limited motion}

FR does not usually take a raw unconstrained step. Instead, it limits the step
size by a temperature parameter $t > 0$:
\[
\vect{x}_i \leftarrow \vect{x}_i
+
\frac{\vect{d}_i}{\norm{\vect{d}_i}}
\min\!\bigl(\norm{\vect{d}_i}, t\bigr).
\]

The terms in this update have the following meaning:
\begin{itemize}[leftmargin=1.5em]
\item $\vect{d}_i / \norm{\vect{d}_i}$ gives only the direction of motion,
\item $\norm{\vect{d}_i}$ is the raw proposed movement magnitude,
\item $\min(\norm{\vect{d}_i}, t)$ caps that motion by the current temperature,
\item $t$ decreases over time according to a cooling schedule.
\end{itemize}

The temperature plays a role similar to simulated annealing:
\begin{itemize}[leftmargin=1.5em]
\item early iterations allow larger moves for global rearrangement,
\item later iterations force smaller moves for stabilization and fine adjustment.
\end{itemize}

\subsection{Algorithmic outline}

A standard FR workflow is:
\begin{enumerate}[leftmargin=1.5em]
\item Initialize vertex positions, often randomly.
\item Set the drawing area $A$, the scale $k$, and an initial temperature $t$.
\item For each iteration:
  \begin{enumerate}[leftmargin=1.5em]
  \item compute repulsive forces between all vertex pairs,
  \item compute attractive forces along edges,
  \item sum them into displacement vectors,
  \item move each vertex by a temperature-limited step,
  \item decrease the temperature.
  \end{enumerate}
\item Stop after a fixed number of iterations or when movements become small.
\end{enumerate}

\subsection{Why the method works}

FR works by enforcing two visual principles:
\begin{itemize}[leftmargin=1.5em]
\item vertices should not collapse onto one another,
\item adjacent vertices should remain close enough to reveal graph connectivity.
\end{itemize}

Because the attractive force acts only on edges, FR is more local than KK in
its structural encoding. The global shape of the layout emerges from repeated
force balancing rather than direct fitting of all graph distances.

\subsection{Strengths and limitations}

\textbf{Strengths}
\begin{itemize}[leftmargin=1.5em]
\item It is simple and intuitive.
\item It often produces aesthetically pleasing drawings with good spacing.
\item It is easy to implement and easy to modify.
\item Many scalable variants exist, including multilevel, grid-based, and
Barnes--Hut style accelerations.
\end{itemize}

\textbf{Limitations}
\begin{itemize}[leftmargin=1.5em]
\item The classical algorithm is still $O(n^2)$ per iteration because of all-pairs repulsion.
\item It does not explicitly preserve graph-theoretic distances.
\item Results can depend strongly on initialization, iteration count, and cooling schedule.
\item Dense graphs or graphs with strong geometric regularity may require careful tuning.
\end{itemize}

\subsection{Computational cost}

For the classical implementation:
\begin{itemize}[leftmargin=1.5em]
\item repulsion over all unordered pairs costs $O(n^2)$ per iteration,
\item attraction over edges costs $O(m)$ per iteration,
\item the total is typically $O(n^2 + m)$ per iteration.
\end{itemize}

Unlike KK, FR does not need all-pairs shortest paths, which is an important
practical advantage.

\section{Similarities between KK and FR}

Although the two methods differ in important ways, they also share a common
force-directed philosophy.

\begin{enumerate}[leftmargin=1.5em]
\item \textbf{Both produce geometric embeddings of graphs.} Each algorithm maps
vertices to points in $\R^2$ (or more generally $\R^d$).

\item \textbf{Both are iterative optimization procedures.} Neither method jumps
directly to the final answer. They repeatedly improve a layout.

\item \textbf{Both balance attraction and separation.} In KK, this appears
through spring energies whose minima enforce desired distances. In FR, it
appears explicitly through attractive and repulsive forces.

\item \textbf{Both are nonconvex.} Multiple local minima are possible, so the
final drawing can depend on initialization and tuning.

\item \textbf{Both can be computationally expensive in naive form.} KK and FR
each involve all-pairs interactions in their classical formulations.

\item \textbf{Both aim at readable graph drawings.} They seek layouts with low
clutter, structural clarity, and visually meaningful geometry.
\end{enumerate}

\section{Differences between KK and FR}

\subsection{Conceptual difference}

The most important conceptual difference is:
\begin{itemize}[leftmargin=1.5em]
\item \textbf{KK is distance-matching.} It tries to realize graph shortest-path
distances in Euclidean space.
\item \textbf{FR is force-balancing.} It tries to balance generic attraction and
repulsion to obtain an aesthetically organized picture.
\end{itemize}

\subsection{Mathematical difference}

KK has an explicit global energy
\[
E(X) = \frac{1}{2} \sum_{i<j} k_{ij}(r_{ij} - \ell_{ij})^2,
\]
so its mathematics is very close to weighted stress minimization.

FR, in its original form, is usually presented as a dynamical procedure with
force calculations and a cooling schedule, not as a single explicit scalar
objective minimized by Newton's method.

\subsection{Structural information used}

\begin{itemize}[leftmargin=1.5em]
\item \textbf{KK uses all graph distances $g_{ij}$.} Every pair of vertices
contributes structurally through $\ell_{ij}$ and $k_{ij}$.
\item \textbf{FR uses edges for attraction and all pairs for repulsion.} It does
not directly encode the full shortest-path metric of the graph.
\end{itemize}

\subsection{Update mechanics}

\begin{itemize}[leftmargin=1.5em]
\item \textbf{KK} selects one vertex with large residual and applies a local
Newton solve using first and second derivatives.
\item \textbf{FR} computes displacement vectors for all vertices and moves all of
them each iteration with a temperature cap.
\end{itemize}

\subsection{Scale parameters}

\begin{itemize}[leftmargin=1.5em]
\item In \textbf{KK}, the desired length $\ell_{ij}$ depends on the pair $(i,j)$,
and the stiffness $k_{ij}$ also depends on the pair.
\item In \textbf{FR}, the parameter $k$ is a single global spacing scale.
\end{itemize}

\subsection{Practical behavior}

\begin{itemize}[leftmargin=1.5em]
\item \textbf{KK} often gives very good layouts for smaller graphs when global
distance fidelity matters.
\item \textbf{FR} is often easier to scale and is widely used as a general
purpose graph drawing method, especially when paired with approximations or
multilevel strategies.
\end{itemize}

\subsection{Comparison table}

\begin{center}
\small
\begin{tabular}{>{\raggedright\arraybackslash}p{0.21\textwidth} >{\raggedright\arraybackslash}p{0.33\textwidth} >{\raggedright\arraybackslash}p{0.33\textwidth}}
\toprule
\textbf{Aspect} & \textbf{Kamada--Kawai} & \textbf{Fruchterman--Reingold} \\
\midrule
Primary viewpoint & Weighted spring energy minimization & Iterative force simulation with cooling \\
Main structural signal & All-pairs shortest-path distances & Edge attraction plus all-pairs repulsion \\
Key scale quantities & Pair-specific $\ell_{ij}$ and $k_{ij}$ & Global spacing parameter $k$ \\
Optimization step & Vertex-wise Newton--Raphson updates & Global displacement updates for all vertices \\
Typical strength & Strong distance preservation & Good visual balance and simple general-purpose behavior \\
Typical weakness & Expensive on large graphs & Does not explicitly preserve graph distances \\
Naive cost & All-pairs distances plus $O(n^2)$ interactions & $O(n^2 + m)$ per iteration \\
\bottomrule
\end{tabular}
\end{center}

\section{When each method is a good fit}

\textbf{KK is often a good fit when}
\begin{itemize}[leftmargin=1.5em]
\item the graph is small or medium-sized,
\item graph-theoretic distance preservation is important,
\item one wants a mathematically explicit objective.
\end{itemize}

\textbf{FR is often a good fit when}
\begin{itemize}[leftmargin=1.5em]
\item one wants a simple and classical force-directed layout,
\item visual spacing and symmetry are the primary goals,
\item scalable approximations or multilevel variants are available.
\end{itemize}

\section{Conclusion}

Kamada--Kawai and Fruchterman--Reingold are both foundational force-directed
methods, but they solve different mathematical problems. KK asks for a drawing
whose Euclidean distances reflect shortest-path distances in the graph. FR asks
for a drawing in which attractive and repulsive forces reach a visually balanced
equilibrium. Because of that difference, KK is more distance-faithful and more
explicitly optimization-based, while FR is simpler, more heuristic, and often
easier to adapt to larger-scale drawing pipelines.

In modern graph layout systems, both methods remain important. KK provides a
clean distance-based formulation, and FR provides a practical force-balancing
template that has inspired many scalable descendants.

\begin{thebibliography}{9}

\bibitem{kamada1989}
T. Kamada and S. Kawai.
\newblock An algorithm for drawing general undirected graphs.
\newblock \emph{Information Processing Letters}, 31(1):7--15, 1989.

\bibitem{fr1991}
T. M. J. Fruchterman and E. M. Reingold.
\newblock Graph drawing by force-directed placement.
\newblock \emph{Software: Practice and Experience}, 21(11):1129--1164, 1991.

\bibitem{eades1984}
P. Eades.
\newblock A heuristic for graph drawing.
\newblock \emph{Congressus Numerantium}, 42:149--160, 1984.

\end{thebibliography}

\end{document}
